% S4: chi duoc \input boi paper2_hujos/supplementary.tex (kiem tra 2026-10-09).
% Truoc day la \section{...} danh so tu dong -> render trong supplementary thanh
% "1 Subsidence Invariance" trong khi than bai trich dan "S4" (peer review F5).
\section*{S4. Subsidence Invariance}

AWD compliance is usually stated as an absolute water-table or ponding depth,
for example $15$\,cm below the surface, the reflooding threshold of the Safe-AWD
treatment in~\cite{balaine2019}. In a subsiding delta
the surface elevation $z_s(t)$ falls at rate $\dot z$, so an absolute threshold
measured against a fixed benchmark drifts.

\begin{proposition}[Drift of an absolute threshold]
\label{prop:subsid}
Over a crediting period of $T_c$ years an absolute threshold expressed against
a fixed datum shifts by $\dot z\,T_c$. For $\dot z\in[1.1,2.5]$\,cm/yr and
$T_c=5$\,yr the shift is $5.5$--$12.5$\,cm, that is $37$--$83\%$ of a
$15$\,cm threshold. A criterion written in the depletion fraction
\eqref{eq:depl} has zero drift, because $FC$ and $WP$ are soil properties
referenced to the root zone and not to an absolute elevation.
\end{proposition}

\begin{proof}
The absolute criterion compares a measured depth to a constant benchmark; if
the benchmark and the surface move apart at $\dot z$, the effective threshold
changes by $\dot z T_c$ over $T_c$ years. The depletion criterion compares
$FC-w_t$ to $(FC-WP)$, both of which are defined within the root zone and
translate rigidly with the surface, so their ratio is invariant to $z_s$.
\end{proof}

The subsidence rates are those measured for the delta by
\cite{minderhoud2017}. The practical consequence is that a scheme anchored on
depletion fraction needs no re-surveying of field benchmarks during a five-year
crediting period, whereas an absolute-depth scheme does.
